\documentclass[letterpaper]{article}
\usepackage[submission]{aaai2027}
\usepackage[hyphens]{url}
\usepackage{graphicx}
\urlstyle{rm}
\def\UrlFont{\rm}
\usepackage{natbib}
\usepackage{caption}
\frenchspacing
\pdfinfo{/TemplateVersion (2027.1)}
% Results: generated by scripts/paper_numbers.py from results/. Read-only.
% GENERATED by scripts/paper_numbers.py from results/. Do not edit by hand.

% --- T behavioural screen, results/t_screen_qwen25vl3b.json (D41)
\newcommand{\nObjects}{469}  % t_screen recomputed n_objects
\newcommand{\nCorrectNeutral}{436}  % t_screen recomputed n_objects_correct_neutral
\newcommand{\nScreenUsable}{286}  % t_screen recomputed n_objects_usable
\newcommand{\screenFoolRate}{65.6}  % t_screen recomputed fool_rate_among_eligible
\newcommand{\screenFoolSmall}{58.3}  % position_vs_size on_object_small fool_rate
\newcommand{\screenFoolMedium}{56.7}  % position_vs_size on_object_medium fool_rate
\newcommand{\screenFoolBeside}{44.2}  % position_vs_size beside_object fool_rate
\newcommand{\pairedOnOnly}{99}  % position_vs_size on_vs_beside.only_on
\newcommand{\pairedBesideOnly}{20}  % position_vs_size on_vs_beside.only_beside
\newcommand{\attackShift}{+8.50}  % t_screen records, mean attack shift
\newcommand{\attackShiftMin}{+2.12}  % t_screen records, min attack shift
\newcommand{\congruentShift}{$-$2.10}  % t_screen records, mean congruent shift
\newcommand{\congruentShiftMax}{+1.38}  % t_screen records, max congruent shift

% --- Replication gate, results/replication_{patching,knockout}_qwen25vl3b.json (D40, D67)
\newcommand{\nLayers}{36}  % replication_patching pc1_scores shape
\newcommand{\nHeadsPerLayer}{16}  % replication_patching pc1_scores shape
\newcommand{\nHeads}{576}  % replication_patching pc1_scores size
\newcommand{\nP}{73}  % replication_patching n_scored
\newcommand{\nPromP}{9}  % replication_patching promoting
\newcommand{\nSupP}{6}  % replication_patching suppressing
\newcommand{\promMatchesLietzow}{identical to}  % promoting == reference_promoting
\newcommand{\repPriorFlip}{89.0}  % knockout grp_promoting prior_flip_rate
\newcommand{\repPriorFlipCount}{65}  % knockout grp_promoting prior_flips
\newcommand{\repVisualFlip}{2.7}  % knockout grp_promoting visual_flip_rate
\newcommand{\lietzowPriorFlip}{84.9}  % knockout reported (Lietzow)
\newcommand{\lietzowVisualFlip}{1.4}  % knockout reported (Lietzow)
\newcommand{\topSingleOne}{L26H5}  % knockout ind_*, ranked by prior_flip_rate
\newcommand{\topSingleOneFlip}{46.6}  % knockout ind_L26H5 prior_flip_rate
\newcommand{\topSingleOneLietzow}{41.1}  % knockout reported (Lietzow)
\newcommand{\topSingleTwo}{L31H3}  % knockout ind_*, ranked by prior_flip_rate
\newcommand{\topSingleTwoFlip}{41.1}  % knockout ind_L31H3 prior_flip_rate
\newcommand{\topSingleTwoLietzow}{37.0}  % knockout reported (Lietzow)

% --- T circuit and J_PT, results/t_patching_qwen25vl3b.json, t_circuit_overlap_qwen25vl3b.json (D70-D72)
\newcommand{\nBoot}{1000}  % parameters.n_boot (same for pg and V overlap)
\newcommand{\nPerm}{10000}  % parameters.n_perm (same for pg and V overlap)
\newcommand{\nT}{244}  % t_patching n_scored
\newcommand{\nPromT}{12}  % t_patching promoting
\newcommand{\nSupT}{6}  % t_patching suppressing
\newcommand{\nShared}{5}  % t_circuit_overlap shared_promoting
\newcommand{\nPonly}{4}  % p_promoting - shared
\newcommand{\nTonly}{7}  % t_promoting - shared
\newcommand{\JPT}{0.312}  % t_circuit_overlap primary.j_pt
\newcommand{\JPTLo}{0.278}  % primary.ci95[0]
\newcommand{\JPTHi}{0.375}  % primary.ci95[1]
\newcommand{\JPTVerdict}{above chance, below self-agreement}  % primary.wording (D70 fixed phrase)
\newcommand{\chanceMean}{0.010}  % primary.chance_floor_mean
\newcommand{\chanceEdge}{0.050}  % primary.chance_floor_p95_at_observed_sizes
\newcommand{\ceilingMedian}{0.900}  % split_half jaccard_promoting.median
\newcommand{\ceilingLo}{0.900}  % split_half jaccard_promoting.p2_5
\newcommand{\nHalfSplits}{1000}  % split_half n_splits
\newcommand{\TselfJ}{1.000}  % stable_core.t_self_agreement_jaccard.median
\newcommand{\rankRho}{0.41}  % rank_correlation.rho
\newcommand{\rankCeiling}{0.77}  % rank_correlation.p_ceiling_full_length
\newcommand{\nSharedObjects}{40}  % shared_objects.n_objects
\newcommand{\JPTshared}{0.400}  % shared_objects.j_pt
\newcommand{\JPTsharedLo}{0.278}  % shared_objects.ci95[0]
\newcommand{\JPTsharedHi}{0.429}  % shared_objects.ci95[1]
\newcommand{\layersPonly}{26--28}  % layers of p_promoting - shared
\newcommand{\layersTonly}{29--33}  % layers of t_promoting - shared
\newcommand{\layersShared}{26--34}  % layers of shared_promoting
\newcommand{\containPinT}{0.56}  % descriptive.containment_p_in_t
\newcommand{\JSup}{0.714}  % descriptive.j_suppressing

% --- H2 cross-ablation, results/cross_ablation_analysis_qwen25vl3b.json (D73)
\newcommand{\hTwoDraws}{100}  % primary P->T n_draws
\newcommand{\hTwoAlpha}{0.025}  % alpha (same in pg_cross_ablation_analysis)
\newcommand{\hTwoMinP}{0.0099}  % min_attainable_p, every zero-ablation test
\newcommand{\hTwoPtoTObs}{2.5}  % cross_ablation P->T observed
\newcommand{\hTwoPtoTLo}{0.8}  % cross_ablation P->T ci95[0]
\newcommand{\hTwoPtoTHi}{4.5}  % cross_ablation P->T ci95[1]
\newcommand{\hTwoPtoTNull}{4.3}  % cross_ablation P->T null_mean
\newcommand{\hTwoPtoTNullEdge}{9.8}  % cross_ablation P->T null_p95
\newcommand{\hTwoPtoTP}{0.75}  % cross_ablation P->T p
\newcommand{\hTwoPtoTCount}{6}  % cross_ablation P->T count
\newcommand{\hTwoPtoTN}{244}  % cross_ablation P->T n
\newcommand{\hTwoTtoPObs}{2.7}  % cross_ablation T->P observed
\newcommand{\hTwoTtoPLo}{0.0}  % cross_ablation T->P ci95[0]
\newcommand{\hTwoTtoPHi}{6.8}  % cross_ablation T->P ci95[1]
\newcommand{\hTwoTtoPNull}{4.9}  % cross_ablation T->P null_mean
\newcommand{\hTwoTtoPNullEdge}{17.8}  % cross_ablation T->P null_p95
\newcommand{\hTwoTtoPP}{0.50}  % cross_ablation T->P p
\newcommand{\hTwoTtoPCount}{2}  % cross_ablation T->P count
\newcommand{\hTwoTtoPN}{73}  % cross_ablation T->P n
\newcommand{\hTwoPtoPObs}{87.7}  % cross_ablation P->P observed
\newcommand{\hTwoPtoPLo}{79.5}  % cross_ablation P->P ci95[0]
\newcommand{\hTwoPtoPHi}{94.5}  % cross_ablation P->P ci95[1]
\newcommand{\hTwoPtoPNull}{7.3}  % cross_ablation P->P null_mean
\newcommand{\hTwoPtoPNullEdge}{20.7}  % cross_ablation P->P null_p95
\newcommand{\hTwoPtoPP}{0.0099}  % cross_ablation P->P p
\newcommand{\hTwoPtoPCount}{64}  % cross_ablation P->P count
\newcommand{\hTwoPtoPN}{73}  % cross_ablation P->P n
\newcommand{\hTwoTtoTObs}{67.2}  % cross_ablation T->T observed
\newcommand{\hTwoTtoTLo}{61.1}  % cross_ablation T->T ci95[0]
\newcommand{\hTwoTtoTHi}{72.5}  % cross_ablation T->T ci95[1]
\newcommand{\hTwoTtoTNull}{5.0}  % cross_ablation T->T null_mean
\newcommand{\hTwoTtoTNullEdge}{10.7}  % cross_ablation T->T null_p95
\newcommand{\hTwoTtoTP}{0.0099}  % cross_ablation T->T p
\newcommand{\hTwoTtoTCount}{164}  % cross_ablation T->T count
\newcommand{\hTwoTtoTN}{244}  % cross_ablation T->T n
\newcommand{\hTwoStoPObs}{83.6}  % cross_ablation shared->P observed
\newcommand{\hTwoStoPLo}{75.3}  % cross_ablation shared->P ci95[0]
\newcommand{\hTwoStoPHi}{91.8}  % cross_ablation shared->P ci95[1]
\newcommand{\hTwoStoPNull}{4.2}  % cross_ablation shared->P null_mean
\newcommand{\hTwoStoPNullEdge}{19.2}  % cross_ablation shared->P null_p95
\newcommand{\hTwoStoPP}{0.0099}  % cross_ablation shared->P p
\newcommand{\hTwoStoPCount}{61}  % cross_ablation shared->P count
\newcommand{\hTwoStoPN}{73}  % cross_ablation shared->P n
\newcommand{\hTwoStoTObs}{0.0}  % cross_ablation shared->T observed
\newcommand{\hTwoStoTLo}{0.0}  % cross_ablation shared->T ci95[0]
\newcommand{\hTwoStoTHi}{0.0}  % cross_ablation shared->T ci95[1]
\newcommand{\hTwoStoTNull}{2.9}  % cross_ablation shared->T null_mean
\newcommand{\hTwoStoTNullEdge}{9.0}  % cross_ablation shared->T null_p95
\newcommand{\hTwoStoTP}{1.00}  % cross_ablation shared->T p
\newcommand{\hTwoStoTCount}{0}  % cross_ablation shared->T count
\newcommand{\hTwoStoTN}{244}  % cross_ablation shared->T n
\newcommand{\hTwoPOtoTObs}{7.8}  % cross_ablation P_only->T observed
\newcommand{\hTwoPOtoTLo}{4.5}  % cross_ablation P_only->T ci95[0]
\newcommand{\hTwoPOtoTHi}{11.5}  % cross_ablation P_only->T ci95[1]
\newcommand{\hTwoPOtoTNull}{2.6}  % cross_ablation P_only->T null_mean
\newcommand{\hTwoPOtoTNullEdge}{6.2}  % cross_ablation P_only->T null_p95
\newcommand{\hTwoPOtoTP}{0.0495}  % cross_ablation P_only->T p
\newcommand{\hTwoPOtoTCount}{19}  % cross_ablation P_only->T count
\newcommand{\hTwoPOtoTN}{244}  % cross_ablation P_only->T n
\newcommand{\hTwoTOtoPObs}{0.0}  % cross_ablation T_only->P observed
\newcommand{\hTwoTOtoPLo}{0.0}  % cross_ablation T_only->P ci95[0]
\newcommand{\hTwoTOtoPHi}{0.0}  % cross_ablation T_only->P ci95[1]
\newcommand{\hTwoTOtoPNull}{2.5}  % cross_ablation T_only->P null_mean
\newcommand{\hTwoTOtoPNullEdge}{9.6}  % cross_ablation T_only->P null_p95
\newcommand{\hTwoTOtoPP}{1.00}  % cross_ablation T_only->P p
\newcommand{\hTwoTOtoPCount}{0}  % cross_ablation T_only->P count
\newcommand{\hTwoTOtoPN}{73}  % cross_ablation T_only->P n
\newcommand{\hTwoTOtoTObs}{70.5}  % own_condition_subsets T_only->T
\newcommand{\hTwoPOtoPObs}{17.8}  % own_condition_subsets P_only->P
\newcommand{\hTwoTOonPShift}{$-$3.74}  % decomposition T_only->P logit_diff_change.mean
\newcommand{\hTwoStoPShift}{+3.62}  % decomposition shared->P logit_diff_change.mean
\newcommand{\hTwoCollatT}{47.9}  % collateral_target_flips P.T
\newcommand{\hTwoCollatTO}{8.2}  % collateral_target_flips P.T_only
\newcommand{\hTwoCleanT}{83.6}  % clean_accuracy accuracy_kept.T
\newcommand{\hTwoCleanTO}{98.5}  % clean_accuracy accuracy_kept.T_only
\newcommand{\hTwoCleanRandom}{99.0}  % clean_accuracy random_mean_rT
\newcommand{\hTwoMeanSentence}{Mean ablation gives the same verdicts.}  % mean_ablation.same_wording_as_zero

% --- H3 mechanism, results/h3_analysis_qwen25vl3b.json (D74)
\newcommand{\hThreeNT}{223}  % n_objects.T
\newcommand{\hThreeNP}{70}  % n_objects.P
\newcommand{\hThreeDropT}{21}  % n_dropped.T
\newcommand{\hThreeDropP}{3}  % n_dropped.P
\newcommand{\hThreeMechanism}{the T-only heads read the note and write its colour}  % mechanism (D74 fixed phrase)
\newcommand{\hThreeAObs}{0.281}  % h3 A observed
\newcommand{\hThreeALo}{0.274}  % h3 A ci95[0]
\newcommand{\hThreeAHi}{0.288}  % h3 A ci95[1]
\newcommand{\hThreeANull}{0.075}  % h3 A null_mean
\newcommand{\hThreeANullEdge}{0.118}  % h3 A null_p95
\newcommand{\hThreeAP}{0.0099}  % h3 A p
\newcommand{\hThreeWObs}{+4.02}  % h3 W observed
\newcommand{\hThreeWLo}{+3.81}  % h3 W ci95[0]
\newcommand{\hThreeWHi}{+4.23}  % h3 W ci95[1]
\newcommand{\hThreeWNull}{+0.02}  % h3 W null_mean
\newcommand{\hThreeWNullEdge}{+0.41}  % h3 W null_p95
\newcommand{\hThreeWP}{0.0099}  % h3 W p
\newcommand{\hThreeASObs}{0.086}  % h3 A:shared observed
\newcommand{\hThreeASLo}{0.083}  % h3 A:shared ci95[0]
\newcommand{\hThreeASHi}{0.090}  % h3 A:shared ci95[1]
\newcommand{\hThreeASNull}{0.068}  % h3 A:shared null_mean
\newcommand{\hThreeASNullEdge}{0.158}  % h3 A:shared null_p95
\newcommand{\hThreeASP}{0.27}  % h3 A:shared p
\newcommand{\hThreeWSObs}{$-$0.18}  % h3 W:shared observed
\newcommand{\hThreeWSLo}{$-$0.29}  % h3 W:shared ci95[0]
\newcommand{\hThreeWSHi}{$-$0.08}  % h3 W:shared ci95[1]
\newcommand{\hThreeWSNull}{+0.02}  % h3 W:shared null_mean
\newcommand{\hThreeWSNullEdge}{+0.16}  % h3 W:shared null_p95
\newcommand{\hThreeWSP}{1.00}  % h3 W:shared p
\newcommand{\hThreeAPOObs}{0.122}  % h3 A:P_only observed
\newcommand{\hThreeAPOLo}{0.117}  % h3 A:P_only ci95[0]
\newcommand{\hThreeAPOHi}{0.127}  % h3 A:P_only ci95[1]
\newcommand{\hThreeAPONull}{0.056}  % h3 A:P_only null_mean
\newcommand{\hThreeAPONullEdge}{0.099}  % h3 A:P_only null_p95
\newcommand{\hThreeAPOP}{0.0099}  % h3 A:P_only p
\newcommand{\hThreeWPOObs}{$-$0.03}  % h3 W:P_only observed
\newcommand{\hThreeWPOLo}{$-$0.05}  % h3 W:P_only ci95[0]
\newcommand{\hThreeWPOHi}{$-$0.02}  % h3 W:P_only ci95[1]
\newcommand{\hThreeWPONull}{$-$0.01}  % h3 W:P_only null_mean
\newcommand{\hThreeWPONullEdge}{+0.03}  % h3 W:P_only null_p95
\newcommand{\hThreeWPOP}{0.77}  % h3 W:P_only p
\newcommand{\hThreeRiseObs}{0.110}  % h3 A_rise:T_only observed
\newcommand{\hThreeRiseLo}{0.100}  % h3 A_rise:T_only ci95[0]
\newcommand{\hThreeRiseHi}{0.119}  % h3 A_rise:T_only ci95[1]
\newcommand{\hThreeRiseNull}{0.026}  % h3 A_rise:T_only null_mean
\newcommand{\hThreeRiseNullEdge}{0.062}  % h3 A_rise:T_only null_p95
\newcommand{\hThreeRiseP}{0.0099}  % h3 A_rise:T_only p
\newcommand{\hThreeRiseSObs}{0.049}  % h3 A_rise:shared observed
\newcommand{\hThreeRiseSLo}{0.045}  % h3 A_rise:shared ci95[0]
\newcommand{\hThreeRiseSHi}{0.053}  % h3 A_rise:shared ci95[1]
\newcommand{\hThreeRiseSNull}{0.039}  % h3 A_rise:shared null_mean
\newcommand{\hThreeRiseSNullEdge}{0.109}  % h3 A_rise:shared null_p95
\newcommand{\hThreeRiseSP}{0.29}  % h3 A_rise:shared p
\newcommand{\hThreeObjCleanObs}{0.256}  % h3 object:clean observed
\newcommand{\hThreeObjCleanLo}{0.245}  % h3 object:clean ci95[0]
\newcommand{\hThreeObjCleanHi}{0.267}  % h3 object:clean ci95[1]
\newcommand{\hThreeObjCleanNull}{0.095}  % h3 object:clean null_mean
\newcommand{\hThreeObjCleanNullEdge}{0.140}  % h3 object:clean null_p95
\newcommand{\hThreeObjCleanP}{0.0099}  % h3 object:clean p
\newcommand{\hThreeColCleanObs}{$-$3.12}  % h3 image_colour:clean observed
\newcommand{\hThreeColCleanLo}{$-$3.27}  % h3 image_colour:clean ci95[0]
\newcommand{\hThreeColCleanHi}{$-$2.97}  % h3 image_colour:clean ci95[1]
\newcommand{\hThreeColCleanNull}{$-$0.06}  % h3 image_colour:clean null_mean
\newcommand{\hThreeColCleanNullEdge}{$-$0.33}  % h3 image_colour:clean null_p5
\newcommand{\hThreeColCleanP}{0.0099}  % h3 image_colour:clean p
\newcommand{\hThreeObjPVObs}{0.346}  % h3 object:P_visual observed
\newcommand{\hThreeObjPVLo}{0.328}  % h3 object:P_visual ci95[0]
\newcommand{\hThreeObjPVHi}{0.365}  % h3 object:P_visual ci95[1]
\newcommand{\hThreeObjPVNull}{0.130}  % h3 object:P_visual null_mean
\newcommand{\hThreeObjPVNullEdge}{0.191}  % h3 object:P_visual null_p95
\newcommand{\hThreeObjPVP}{0.0099}  % h3 object:P_visual p
\newcommand{\hThreeColPVObs}{$-$4.27}  % h3 image_colour:P_visual observed
\newcommand{\hThreeColPVLo}{$-$4.65}  % h3 image_colour:P_visual ci95[0]
\newcommand{\hThreeColPVHi}{$-$3.92}  % h3 image_colour:P_visual ci95[1]
\newcommand{\hThreeColPVNull}{+0.03}  % h3 image_colour:P_visual null_mean
\newcommand{\hThreeColPVNullEdge}{$-$0.30}  % h3 image_colour:P_visual null_p5
\newcommand{\hThreeColPVP}{0.0099}  % h3 image_colour:P_visual p
\newcommand{\hThreeRoutingObs}{0.230}  % h3 routing:P observed
\newcommand{\hThreeRoutingLo}{0.214}  % h3 routing:P ci95[0]
\newcommand{\hThreeRoutingHi}{0.246}  % h3 routing:P ci95[1]
\newcommand{\hThreeRoutingNull}{0.087}  % h3 routing:P null_mean
\newcommand{\hThreeRoutingNullEdge}{0.171}  % h3 routing:P null_p95
\newcommand{\hThreeRoutingP}{0.0198}  % h3 routing:P p
\newcommand{\noteShare}{3.1}  % descriptive.note_token_share
\newcommand{\enrichInc}{13.2}  % note_enrichment T_only incongruent
\newcommand{\enrichNeu}{9.0}  % note_enrichment T_only neutral
\newcommand{\foolSmall}{100}  % spatial on_object_small fool_rate
\newcommand{\foolMedium}{87}  % spatial on_object_medium fool_rate
\newcommand{\foolBeside}{68}  % spatial beside_object fool_rate

% --- V condition, results/v_patching_qwen25vl3b.json, v_circuit_overlap_qwen25vl3b.json (D75)
\newcommand{\nVselected}{392}  % v_patching n_selected
\newcommand{\nV}{244}  % v_patching n_scored
\newcommand{\nPromV}{7}  % v_circuit_overlap v_promoting
\newcommand{\nSupV}{1}  % v_circuit_overlap v_suppressing
\newcommand{\layersV}{31--34}  % layers of v_promoting
\newcommand{\VceilingLo}{0.857}  % v_ceiling jaccard_promoting.p2_5
\newcommand{\JVT}{0.583}  % j_vt primary.j_pt
\newcommand{\JVTLo}{0.462}  % j_vt primary.ci95[0]
\newcommand{\JVTHi}{0.583}  % j_vt primary.ci95[1]
\newcommand{\JVTVerdict}{above chance, below self-agreement}  % j_vt primary.wording (D70 fixed phrase)
\newcommand{\chanceEdgeVT}{0.056}  % j_vt primary.chance_floor_p95_at_observed_sizes
\newcommand{\ceilingVTLo}{0.857}  % j_vt ceiling_binding V, p2_5
\newcommand{\JPV}{0.231}  % j_pv primary.j_pt
\newcommand{\JPVLo}{0.214}  % j_pv primary.ci95[0]
\newcommand{\JPVHi}{0.250}  % j_pv primary.ci95[1]
\newcommand{\JPVVerdict}{above chance, below self-agreement}  % j_pv primary.wording (D70 fixed phrase)
\newcommand{\chanceEdgePV}{0.067}  % j_pv primary.chance_floor_p95_at_observed_sizes
\newcommand{\ceilingPVLo}{0.857}  % j_pv ceiling_binding V, p2_5
\newcommand{\nSharedPV}{3}  % j_pv shared_promoting
\newcommand{\dJ}{+0.271}  % delta_j.delta_j
\newcommand{\dJLo}{+0.163}  % delta_j.ci95[0]
\newcommand{\dJHi}{+0.289}  % delta_j.ci95[1]
\newcommand{\dJVerdict}{$J_{VT} > J_{PT}$}  % delta_j.wording (D75 fixed phrase)
\newcommand{\nTonlyInV}{4}  % descriptive.t_only_heads_in_v.n
\newcommand{\nSharedInV}{3}  % descriptive.shared_heads_in_v.n
\newcommand{\containVinP}{0.43}  % descriptive.containment_v_in_p_promoting
\newcommand{\nSharedObjectsVT}{143}  % j_vt shared_objects.n_objects
\newcommand{\JVTshared}{0.500}  % j_vt shared_objects.j_pt
\newcommand{\JVTsharedLo}{0.385}  % j_vt shared_objects.ci95[0]
\newcommand{\JVTsharedHi}{0.583}  % j_vt shared_objects.ci95[1]

% --- PaliGemma replication, results/pg_replication_{patching,knockout}_paligemma2_3b.json, split_half_ceiling_p_paligemma2_3b.json (D76, D76a)
\newcommand{\pgLayers}{26}  % pg_replication_patching pc1_scores shape
\newcommand{\pgHeadsPerLayer}{8}  % pg_replication_patching pc1_scores shape
\newcommand{\pgHeads}{208}  % pg_replication_patching pc1_scores size
\newcommand{\pgNP}{121}  % pg_replication_patching n_scored
\newcommand{\pgNPromP}{8}  % pg_replication_patching promoting
\newcommand{\pgNSupP}{2}  % pg_replication_patching suppressing
\newcommand{\pgMatchesLietzow}{identical to}  % promoting and suppressing == reference
\newcommand{\pgPriorFlip}{95.9}  % pg knockout grp_promoting prior_flip_rate
\newcommand{\pgPriorFlipCount}{116}  % pg knockout grp_promoting prior_flips
\newcommand{\pgVisualFlip}{1.7}  % pg knockout grp_promoting visual_flip_rate
\newcommand{\pgLietzowPriorFlip}{95.9}  % pg knockout reported (Lietzow)
\newcommand{\pgLietzowVisualFlip}{0.8}  % pg knockout reported (Lietzow)
\newcommand{\pgTopOne}{L15H7}  % pg knockout ind_*, ranked by prior_flip_rate
\newcommand{\pgTopOneFlip}{56.2}  % pg knockout ind_L15H7 prior_flip_rate
\newcommand{\pgTopOneLietzow}{58.7}  % pg knockout reported (Lietzow)
\newcommand{\pgTopTwo}{L20H3}  % pg knockout ind_*, ranked by prior_flip_rate
\newcommand{\pgTopTwoFlip}{38.8}  % pg knockout ind_L20H3 prior_flip_rate
\newcommand{\pgTopTwoLietzow}{36.4}  % pg knockout reported (Lietzow)
\newcommand{\pgCeilingMedian}{1.000}  % pg split_half jaccard_promoting.median
\newcommand{\pgCeilingLo}{0.875}  % pg split_half jaccard_promoting.p2_5

% --- PaliGemma T circuit and J_PT, results/pg_t_patching_paligemma2_3b.json, pg_t_circuit_overlap_paligemma2_3b.json (D77)
\newcommand{\pgNT}{302}  % pg_t_patching n_scored
\newcommand{\pgNPromT}{8}  % pg_t_patching promoting
\newcommand{\pgNSupT}{3}  % pg_t_patching suppressing
\newcommand{\pgNShared}{4}  % pg_t_circuit_overlap shared_promoting
\newcommand{\pgJPT}{0.333}  % pg primary.j_pt
\newcommand{\pgJPTLo}{0.333}  % pg primary.ci95[0]
\newcommand{\pgJPTHi}{0.364}  % pg primary.ci95[1]
\newcommand{\pgJPTVerdict}{above chance, below self-agreement}  % pg primary.wording (D70 fixed phrase)
\newcommand{\pgChanceEdge}{0.067}  % pg primary.chance_floor_p95_at_observed_sizes
\newcommand{\pgNSharedObjects}{69}  % pg shared_objects.n_objects
\newcommand{\pgJPTshared}{0.286}  % pg shared_objects.j_pt
\newcommand{\pgJPTsharedLo}{0.267}  % pg shared_objects.ci95[0]
\newcommand{\pgJPTsharedHi}{0.333}  % pg shared_objects.ci95[1]
\newcommand{\pgRankRho}{0.65}  % pg rank_correlation.rho
\newcommand{\pgRankCeiling}{0.95}  % pg rank_correlation.p_ceiling_full_length

% --- PaliGemma H2 cross-ablation, results/pg_cross_ablation_analysis_paligemma2_3b.json (D78)
\newcommand{\pgDraws}{100}  % pg primary P->T n_draws
\newcommand{\pgPtoTObs}{7.9}  % pg_cross_ablation P->T observed
\newcommand{\pgPtoTLo}{5.3}  % pg_cross_ablation P->T ci95[0]
\newcommand{\pgPtoTHi}{11.3}  % pg_cross_ablation P->T ci95[1]
\newcommand{\pgPtoTNull}{3.5}  % pg_cross_ablation P->T null_mean
\newcommand{\pgPtoTNullEdge}{10.6}  % pg_cross_ablation P->T null_p95
\newcommand{\pgPtoTP}{0.12}  % pg_cross_ablation P->T p
\newcommand{\pgPtoTCount}{24}  % pg_cross_ablation P->T count
\newcommand{\pgPtoTN}{302}  % pg_cross_ablation P->T n
\newcommand{\pgTtoPObs}{28.9}  % pg_cross_ablation T->P observed
\newcommand{\pgTtoPLo}{20.7}  % pg_cross_ablation T->P ci95[0]
\newcommand{\pgTtoPHi}{37.2}  % pg_cross_ablation T->P ci95[1]
\newcommand{\pgTtoPNull}{5.6}  % pg_cross_ablation T->P null_mean
\newcommand{\pgTtoPNullEdge}{13.2}  % pg_cross_ablation T->P null_p95
\newcommand{\pgTtoPP}{0.0099}  % pg_cross_ablation T->P p
\newcommand{\pgTtoPCount}{35}  % pg_cross_ablation T->P count
\newcommand{\pgTtoPN}{121}  % pg_cross_ablation T->P n
\newcommand{\pgPtoPObs}{95.9}  % pg_cross_ablation P->P observed
\newcommand{\pgPtoPLo}{92.6}  % pg_cross_ablation P->P ci95[0]
\newcommand{\pgPtoPHi}{99.2}  % pg_cross_ablation P->P ci95[1]
\newcommand{\pgPtoPNull}{1.9}  % pg_cross_ablation P->P null_mean
\newcommand{\pgPtoPNullEdge}{7.5}  % pg_cross_ablation P->P null_p95
\newcommand{\pgPtoPP}{0.0099}  % pg_cross_ablation P->P p
\newcommand{\pgPtoPCount}{116}  % pg_cross_ablation P->P count
\newcommand{\pgPtoPN}{121}  % pg_cross_ablation P->P n
\newcommand{\pgTtoTObs}{53.6}  % pg_cross_ablation T->T observed
\newcommand{\pgTtoTLo}{48.0}  % pg_cross_ablation T->T ci95[0]
\newcommand{\pgTtoTHi}{58.9}  % pg_cross_ablation T->T ci95[1]
\newcommand{\pgTtoTNull}{2.1}  % pg_cross_ablation T->T null_mean
\newcommand{\pgTtoTNullEdge}{5.0}  % pg_cross_ablation T->T null_p95
\newcommand{\pgTtoTP}{0.0099}  % pg_cross_ablation T->T p
\newcommand{\pgTtoTCount}{162}  % pg_cross_ablation T->T count
\newcommand{\pgTtoTN}{302}  % pg_cross_ablation T->T n
\newcommand{\pgStoPObs}{75.2}  % pg_cross_ablation shared->P observed
\newcommand{\pgStoPLo}{67.8}  % pg_cross_ablation shared->P ci95[0]
\newcommand{\pgStoPHi}{82.6}  % pg_cross_ablation shared->P ci95[1]
\newcommand{\pgStoPNull}{1.8}  % pg_cross_ablation shared->P null_mean
\newcommand{\pgStoPNullEdge}{4.1}  % pg_cross_ablation shared->P null_p95
\newcommand{\pgStoPP}{0.0099}  % pg_cross_ablation shared->P p
\newcommand{\pgStoPCount}{91}  % pg_cross_ablation shared->P count
\newcommand{\pgStoPN}{121}  % pg_cross_ablation shared->P n
\newcommand{\pgStoTObs}{8.9}  % pg_cross_ablation shared->T observed
\newcommand{\pgStoTLo}{6.0}  % pg_cross_ablation shared->T ci95[0]
\newcommand{\pgStoTHi}{12.3}  % pg_cross_ablation shared->T ci95[1]
\newcommand{\pgStoTNull}{1.5}  % pg_cross_ablation shared->T null_mean
\newcommand{\pgStoTNullEdge}{3.7}  % pg_cross_ablation shared->T null_p95
\newcommand{\pgStoTP}{0.0099}  % pg_cross_ablation shared->T p
\newcommand{\pgStoTCount}{27}  % pg_cross_ablation shared->T count
\newcommand{\pgStoTN}{302}  % pg_cross_ablation shared->T n
\newcommand{\pgPOtoTObs}{2.0}  % pg_cross_ablation P_only->T observed
\newcommand{\pgPOtoTLo}{0.7}  % pg_cross_ablation P_only->T ci95[0]
\newcommand{\pgPOtoTHi}{3.6}  % pg_cross_ablation P_only->T ci95[1]
\newcommand{\pgPOtoTNull}{3.6}  % pg_cross_ablation P_only->T null_mean
\newcommand{\pgPOtoTNullEdge}{13.2}  % pg_cross_ablation P_only->T null_p95
\newcommand{\pgPOtoTP}{0.48}  % pg_cross_ablation P_only->T p
\newcommand{\pgPOtoTCount}{6}  % pg_cross_ablation P_only->T count
\newcommand{\pgPOtoTN}{302}  % pg_cross_ablation P_only->T n
\newcommand{\pgTOtoPObs}{0.0}  % pg_cross_ablation T_only->P observed
\newcommand{\pgTOtoPLo}{0.0}  % pg_cross_ablation T_only->P ci95[0]
\newcommand{\pgTOtoPHi}{0.0}  % pg_cross_ablation T_only->P ci95[1]
\newcommand{\pgTOtoPNull}{5.5}  % pg_cross_ablation T_only->P null_mean
\newcommand{\pgTOtoPNullEdge}{10.8}  % pg_cross_ablation T_only->P null_p95
\newcommand{\pgTOtoPP}{1.00}  % pg_cross_ablation T_only->P p
\newcommand{\pgTOtoPCount}{0}  % pg_cross_ablation T_only->P count
\newcommand{\pgTOtoPN}{121}  % pg_cross_ablation T_only->P n
\newcommand{\pgPtoTVerdict}{ablating $\mathcal{H}_P$: not distinguishable from layer-matched random heads at this $n$}  % pg primary P->T wording (D73 fixed phrase)
\newcommand{\pgTtoPVerdict}{ablating $\mathcal{H}_T$ defeats the P override more than layer-matched random heads}  % pg primary T->P wording (D73 fixed phrase)
\newcommand{\pgTOtoTObs}{42.1}  % pg own_condition_subsets T_only->T
\newcommand{\pgPOtoPObs}{71.9}  % pg own_condition_subsets P_only->P
\newcommand{\pgTOonPShift}{$-$1.51}  % pg decomposition T_only->P logit_diff_change.mean
\newcommand{\pgStoPShift}{+1.99}  % pg decomposition shared->P logit_diff_change.mean
\newcommand{\pgMeanDraws}{50}  % pg mean_ablation P->T n_draws
\newcommand{\pgMeanPtoTObs}{11.3}  % pg mean_ablation P->T observed
\newcommand{\pgMeanPtoTLo}{7.6}  % pg mean_ablation P->T ci95[0]
\newcommand{\pgMeanPtoTHi}{15.2}  % pg mean_ablation P->T ci95[1]
\newcommand{\pgMeanPtoTNull}{3.2}  % pg mean_ablation P->T null_mean
\newcommand{\pgMeanPtoTNullEdge}{7.2}  % pg mean_ablation P->T null_p95
\newcommand{\pgMeanPtoTP}{0.0196}  % pg mean_ablation P->T p
\newcommand{\pgMeanPtoTVerdict}{ablating $\mathcal{H}_P$ defeats the T override more than layer-matched random heads}  % pg mean_ablation P->T wording
\newcommand{\pgMeanTtoPObs}{57.0}  % pg mean_ablation T->P observed
\newcommand{\pgMeanTtoPLo}{47.9}  % pg mean_ablation T->P ci95[0]
\newcommand{\pgMeanTtoPHi}{66.1}  % pg mean_ablation T->P ci95[1]
\newcommand{\pgMeanTtoPNull}{1.9}  % pg mean_ablation T->P null_mean
\newcommand{\pgMeanTtoPNullEdge}{5.4}  % pg mean_ablation T->P null_p95
\newcommand{\pgMeanTtoPP}{0.0196}  % pg mean_ablation T->P p
\newcommand{\pgMeanTtoPVerdict}{ablating $\mathcal{H}_T$ defeats the P override more than layer-matched random heads}  % pg mean_ablation T->P wording
\newcommand{\pgMeanSentence}{Mean ablation changes at least one verdict, so the verdict is not robust to the ablation type.}  % pg mean_ablation.same_wording_as_zero

% --- PaliGemma H3 mechanism, results/pg_h3_analysis_paligemma2_3b.json (D80, D80a)
\newcommand{\pgThreeNT}{297}  % pg_h3 n_objects.T
\newcommand{\pgThreeNP}{115}  % pg_h3 n_objects.P
\newcommand{\pgThreeDropT}{5}  % pg_h3 n_dropped.T (drift under eager attention)
\newcommand{\pgThreeDropP}{6}  % pg_h3 n_dropped.P (drift under eager attention)
\newcommand{\pgThreeMechanism}{the T-only heads read the note and write its colour}  % pg_h3 mechanism (D74 fixed phrase)
\newcommand{\pgThreeAObs}{0.150}  % pg_h3 A observed
\newcommand{\pgThreeALo}{0.145}  % pg_h3 A ci95[0]
\newcommand{\pgThreeAHi}{0.154}  % pg_h3 A ci95[1]
\newcommand{\pgThreeANull}{0.037}  % pg_h3 A null_mean
\newcommand{\pgThreeANullEdge}{0.070}  % pg_h3 A null_p95
\newcommand{\pgThreeAP}{0.0099}  % pg_h3 A p
\newcommand{\pgThreeWObs}{+2.70}  % pg_h3 W observed
\newcommand{\pgThreeWLo}{+2.54}  % pg_h3 W ci95[0]
\newcommand{\pgThreeWHi}{+2.85}  % pg_h3 W ci95[1]
\newcommand{\pgThreeWNull}{+0.05}  % pg_h3 W null_mean
\newcommand{\pgThreeWNullEdge}{+0.24}  % pg_h3 W null_p95
\newcommand{\pgThreeWP}{0.0099}  % pg_h3 W p
\newcommand{\pgThreeASObs}{0.037}  % pg_h3 A:shared observed
\newcommand{\pgThreeASLo}{0.036}  % pg_h3 A:shared ci95[0]
\newcommand{\pgThreeASHi}{0.039}  % pg_h3 A:shared ci95[1]
\newcommand{\pgThreeASNull}{0.006}  % pg_h3 A:shared null_mean
\newcommand{\pgThreeASNullEdge}{0.010}  % pg_h3 A:shared null_p95
\newcommand{\pgThreeASP}{0.0099}  % pg_h3 A:shared p
\newcommand{\pgThreeWSObs}{+0.61}  % pg_h3 W:shared observed
\newcommand{\pgThreeWSLo}{+0.49}  % pg_h3 W:shared ci95[0]
\newcommand{\pgThreeWSHi}{+0.72}  % pg_h3 W:shared ci95[1]
\newcommand{\pgThreeWSNull}{$-$0.07}  % pg_h3 W:shared null_mean
\newcommand{\pgThreeWSNullEdge}{+0.04}  % pg_h3 W:shared null_p95
\newcommand{\pgThreeWSP}{0.0099}  % pg_h3 W:shared p
\newcommand{\pgThreeAPOObs}{0.097}  % pg_h3 A:P_only observed
\newcommand{\pgThreeAPOLo}{0.092}  % pg_h3 A:P_only ci95[0]
\newcommand{\pgThreeAPOHi}{0.101}  % pg_h3 A:P_only ci95[1]
\newcommand{\pgThreeAPONull}{0.049}  % pg_h3 A:P_only null_mean
\newcommand{\pgThreeAPONullEdge}{0.109}  % pg_h3 A:P_only null_p95
\newcommand{\pgThreeAPOP}{0.17}  % pg_h3 A:P_only p
\newcommand{\pgThreeWPOObs}{+0.01}  % pg_h3 W:P_only observed
\newcommand{\pgThreeWPOLo}{$-$0.03}  % pg_h3 W:P_only ci95[0]
\newcommand{\pgThreeWPOHi}{+0.05}  % pg_h3 W:P_only ci95[1]
\newcommand{\pgThreeWPONull}{+0.06}  % pg_h3 W:P_only null_mean
\newcommand{\pgThreeWPONullEdge}{+0.20}  % pg_h3 W:P_only null_p95
\newcommand{\pgThreeWPOP}{0.63}  % pg_h3 W:P_only p
\newcommand{\pgThreeRiseObs}{0.025}  % pg_h3 A_rise:T_only observed
\newcommand{\pgThreeRiseLo}{0.021}  % pg_h3 A_rise:T_only ci95[0]
\newcommand{\pgThreeRiseHi}{0.029}  % pg_h3 A_rise:T_only ci95[1]
\newcommand{\pgThreeRiseNull}{$-$0.011}  % pg_h3 A_rise:T_only null_mean
\newcommand{\pgThreeRiseNullEdge}{0.002}  % pg_h3 A_rise:T_only null_p95
\newcommand{\pgThreeRiseP}{0.0099}  % pg_h3 A_rise:T_only p
\newcommand{\pgThreeRiseSObs}{0.007}  % pg_h3 A_rise:shared observed
\newcommand{\pgThreeRiseSLo}{0.006}  % pg_h3 A_rise:shared ci95[0]
\newcommand{\pgThreeRiseSHi}{0.008}  % pg_h3 A_rise:shared ci95[1]
\newcommand{\pgThreeRiseSNull}{$-$0.000}  % pg_h3 A_rise:shared null_mean
\newcommand{\pgThreeRiseSNullEdge}{0.002}  % pg_h3 A_rise:shared null_p95
\newcommand{\pgThreeRiseSP}{0.0099}  % pg_h3 A_rise:shared p
\newcommand{\pgThreeObjCleanObs}{0.234}  % pg_h3 object:clean observed
\newcommand{\pgThreeObjCleanLo}{0.223}  % pg_h3 object:clean ci95[0]
\newcommand{\pgThreeObjCleanHi}{0.245}  % pg_h3 object:clean ci95[1]
\newcommand{\pgThreeObjCleanNull}{0.085}  % pg_h3 object:clean null_mean
\newcommand{\pgThreeObjCleanNullEdge}{0.107}  % pg_h3 object:clean null_p95
\newcommand{\pgThreeObjCleanP}{0.0099}  % pg_h3 object:clean p
\newcommand{\pgThreeColCleanObs}{$-$2.84}  % pg_h3 image_colour:clean observed
\newcommand{\pgThreeColCleanLo}{$-$2.94}  % pg_h3 image_colour:clean ci95[0]
\newcommand{\pgThreeColCleanHi}{$-$2.74}  % pg_h3 image_colour:clean ci95[1]
\newcommand{\pgThreeColCleanNull}{+0.01}  % pg_h3 image_colour:clean null_mean
\newcommand{\pgThreeColCleanNullEdge}{$-$0.09}  % pg_h3 image_colour:clean null_p5
\newcommand{\pgThreeColCleanP}{0.0099}  % pg_h3 image_colour:clean p
\newcommand{\pgThreeObjPVObs}{0.254}  % pg_h3 object:P_visual observed
\newcommand{\pgThreeObjPVLo}{0.240}  % pg_h3 object:P_visual ci95[0]
\newcommand{\pgThreeObjPVHi}{0.271}  % pg_h3 object:P_visual ci95[1]
\newcommand{\pgThreeObjPVNull}{0.097}  % pg_h3 object:P_visual null_mean
\newcommand{\pgThreeObjPVNullEdge}{0.121}  % pg_h3 object:P_visual null_p95
\newcommand{\pgThreeObjPVP}{0.0099}  % pg_h3 object:P_visual p
\newcommand{\pgThreeColPVObs}{$-$3.04}  % pg_h3 image_colour:P_visual observed
\newcommand{\pgThreeColPVLo}{$-$3.23}  % pg_h3 image_colour:P_visual ci95[0]
\newcommand{\pgThreeColPVHi}{$-$2.83}  % pg_h3 image_colour:P_visual ci95[1]
\newcommand{\pgThreeColPVNull}{+0.04}  % pg_h3 image_colour:P_visual null_mean
\newcommand{\pgThreeColPVNullEdge}{$-$0.10}  % pg_h3 image_colour:P_visual null_p5
\newcommand{\pgThreeColPVP}{0.0099}  % pg_h3 image_colour:P_visual p
\newcommand{\pgThreeRoutingObs}{0.058}  % pg_h3 routing:P observed
\newcommand{\pgThreeRoutingLo}{0.050}  % pg_h3 routing:P ci95[0]
\newcommand{\pgThreeRoutingHi}{0.067}  % pg_h3 routing:P ci95[1]
\newcommand{\pgThreeRoutingNull}{0.018}  % pg_h3 routing:P null_mean
\newcommand{\pgThreeRoutingNullEdge}{0.059}  % pg_h3 routing:P null_p95
\newcommand{\pgThreeRoutingP}{0.0792}  % pg_h3 routing:P p
\newcommand{\pgNoteShare}{3.0}  % pg_h3 descriptive.note_token_share
\newcommand{\pgEnrichInc}{7.4}  % pg_h3 note_enrichment T_only incongruent
\newcommand{\pgEnrichNeu}{6.1}  % pg_h3 note_enrichment T_only neutral
\newcommand{\pgFoolSmall}{100.0}  % pg_h3 spatial on_object_small fool_rate
\newcommand{\pgFoolMedium}{96.6}  % pg_h3 spatial on_object_medium fool_rate
\newcommand{\pgFoolBeside}{71.4}  % pg_h3 spatial beside_object fool_rate

% --- PaliGemma V condition, results/pg_v_patching_paligemma2_3b.json, pg_v_circuit_overlap_paligemma2_3b.json (D79, D79a, D79b)
\newcommand{\pgNVselected}{380}  % pg_v_patching n_selected
\newcommand{\pgNV}{302}  % pg_v_patching n_scored
\newcommand{\pgVcap}{302}  % pg_v_patching cap (D79a)
\newcommand{\pgNPromV}{10}  % pg_v_circuit_overlap v_promoting
\newcommand{\pgNSupV}{2}  % pg_v_circuit_overlap v_suppressing
\newcommand{\pgLayersV}{19--25}  % layers of pg v_promoting
\newcommand{\pgVceilingLo}{1.000}  % pg v_ceiling jaccard_promoting.p2_5
\newcommand{\pgJVT}{0.800}  % pg j_vt primary.j_pt
\newcommand{\pgJVTLo}{0.700}  % pg j_vt primary.ci95[0]
\newcommand{\pgJVTHi}{0.800}  % pg j_vt primary.ci95[1]
\newcommand{\pgJVTVerdict}{above chance, below self-agreement}  % pg j_vt primary.wording (D70 fixed phrase)
\newcommand{\pgJVTChanceEdge}{0.059}  % pg j_vt primary.chance_floor_p95_at_observed_sizes
\newcommand{\pgCeilingVTLo}{0.875}  % pg j_vt ceiling_binding T, p2_5
\newcommand{\pgJPV}{0.385}  % pg j_pv primary.j_pt
\newcommand{\pgJPVLo}{0.308}  % pg j_pv primary.ci95[0]
\newcommand{\pgJPVHi}{0.385}  % pg j_pv primary.ci95[1]
\newcommand{\pgJPVVerdict}{above chance, below self-agreement}  % pg j_pv primary.wording (D70 fixed phrase)
\newcommand{\pgJPVChanceEdge}{0.059}  % pg j_pv primary.chance_floor_p95_at_observed_sizes
\newcommand{\pgCeilingPVLo}{0.875}  % pg j_pv ceiling_binding P, p2_5
\newcommand{\pgNSharedPV}{5}  % pg j_pv shared_promoting
\newcommand{\pgNSharedVT}{8}  % pg j_vt shared_promoting
\newcommand{\pgdJ}{+0.467}  % pg delta_j.delta_j
\newcommand{\pgdJLo}{+0.336}  % pg delta_j.ci95[0]
\newcommand{\pgdJHi}{+0.467}  % pg delta_j.ci95[1]
\newcommand{\pgdJVerdict}{$J_{VT} > J_{PT}$}  % pg delta_j.wording (D75 fixed phrase)
\newcommand{\pgNTonlyInV}{4}  % pg descriptive.t_only_heads_in_v.n
\newcommand{\pgNTonly}{4}  % pg descriptive.t_only_heads_in_v.of
\newcommand{\pgNSharedInV}{4}  % pg descriptive.shared_heads_in_v.n
\newcommand{\pgContainVinP}{0.50}  % pg descriptive.containment_v_in_p_promoting
\newcommand{\pgNSharedObjectsVT}{198}  % pg j_vt shared_objects.n_objects
\newcommand{\pgJVTshared}{0.800}  % pg j_vt shared_objects.j_pt
\newcommand{\pgJVTsharedLo}{0.700}  % pg j_vt shared_objects.ci95[0]
\newcommand{\pgJVTsharedHi}{0.800}  % pg j_vt shared_objects.ci95[1]
\newcommand{\pgSoftcapSentence}{Analysing the logits as recorded gives the same head sets and verdict.}  % pg d79b.same_as_primary

% --- Head lists, the promoting/suppressing fields of each discovery file (appendix table)
\newcommand{\listPProm}{L26H0, L26H5, L26H6, L27H1, L27H4, L28H3, L31H3, L31H7, L34H14}  % replication_patching_qwen25vl3b.json promoting
\newcommand{\listPSup}{L26H3, L27H2, L28H4, L30H3, L31H8, L34H10}  % replication_patching_qwen25vl3b.json suppressing
\newcommand{\listTProm}{L26H0, L26H5, L29H2, L31H3, L31H7, L32H6, L32H7, L32H13, L33H1, L33H3, L33H14, L34H14}  % t_patching_qwen25vl3b.json promoting
\newcommand{\listTSup}{L26H3, L26H7, L28H4, L30H3, L31H8, L34H10}  % t_patching_qwen25vl3b.json suppressing
\newcommand{\listVProm}{L31H3, L31H7, L32H6, L32H13, L33H3, L33H14, L34H14}  % v_patching_qwen25vl3b.json promoting
\newcommand{\listVSup}{L34H10}  % v_patching_qwen25vl3b.json suppressing
\newcommand{\pgListPProm}{L15H7, L16H4, L17H7, L20H3, L22H1, L24H3, L25H2, L25H4}  % pg_replication_patching_paligemma2_3b.json promoting
\newcommand{\pgListPSup}{L17H4, L25H5}  % pg_replication_patching_paligemma2_3b.json suppressing
\newcommand{\pgListTProm}{L20H3, L20H7, L22H4, L22H5, L23H7, L24H3, L25H2, L25H4}  % pg_t_patching_paligemma2_3b.json promoting
\newcommand{\pgListTSup}{L17H4, L24H2, L25H5}  % pg_t_patching_paligemma2_3b.json suppressing
\newcommand{\pgListVProm}{L19H1, L20H3, L20H7, L22H1, L22H4, L22H5, L23H7, L24H3, L25H2, L25H4}  % pg_v_circuit_overlap_paligemma2_3b.json v_promoting
\newcommand{\pgListVSup}{L25H3, L25H5}  % pg_v_circuit_overlap_paligemma2_3b.json v_suppressing


% Notation
\newcommand{\HP}{\mathcal{H}_P}
\newcommand{\HT}{\mathcal{H}_T}
\newcommand{\HV}{\mathcal{H}_V}
\newcommand{\cco}{c_{o}}
\newcommand{\ccf}{c_{cf}}

% Drafting only: both must be gone before submission.
\newcommand{\todo}[1]{[TODO: #1]}
\newcommand{\figplaceholder}[3]{\fbox{\parbox[c][#2][c]{#1}{\centering
  [Figure placeholder: #3]}}}
\usepackage{amsmath}
\usepackage{booktabs}
\setcounter{secnumdepth}{0}

\title{Does In-Image Text Hijack the Same Heads as Prompted Priors?\\
A Circuit-Overlap Test in Two Vision-Language Models}
\author{Anonymous Submission}
\affiliations{}

\begin{document}
\maketitle

\begin{abstract}
Vision-language models answer from the image by default, but other signals can override it. Prior work found sparse late attention heads let prompts asking an object's \emph{usual} colour override recoloured images. We ask whether the same heads let a colour word \emph{written in the image}, a typographic attack, override the object's appearance. Across Qwen2.5-VL-3B and PaliGemma2-3B, we hold task, objects and circuit-discovery pipeline fixed, varying only the competing signal's source. Our pipeline re-discovers the published prior-override circuit exactly in both models. The typographic circuit partly overlaps it (Jaccard \JPT{} and \pgJPT): above chance but well below the prior circuit's self-agreement. In Qwen, overlap is in selection, not cause: ablating prior-override heads leaves the typographic attack intact. Instead, later heads attend to the note and write its colour; without a note, they attend to the recoloured object and write its colour. The image-varying circuit overlaps T more than P does (\JVT{} vs \JPT). In PaliGemma, causal separation is partial and depends on ablation type. In-image text appears to exploit a pathway moving colour from image to answer, not the circuit letting prompted knowledge override vision.
\end{abstract}

\section{Introduction}

When a vision-language model (VLM) is asked about an image, the image is not the
only possible source of the answer. The model has memorised what objects usually
look like, and text can appear both in the prompt and inside the image. How a
model arbitrates between visual evidence and these competing signals matters for
reliability: an answer that silently follows the wrong source is a failure that
looks like a success.

\citet{lietzow2026} studied one such arbitration with a recoloured-object
benchmark~\citep{golovanevsky2025counterfact}. Shown a blue strawberry and asked
what colour strawberries \emph{usually} are, a model must ignore the image and
answer from memory. Models mostly fail, following the image; visual grounding is
the default. When the override does succeed, it depends on a small set of
attention heads in the second half of the network. In Qwen-VL these heads shift
attention from image tokens toward text tokens, and ablating them as a group
restores the visual answer~\citep{lietzow2026}.

Typographic attacks put the competing signal somewhere else: \emph{inside} the
image~\citep{goh2021multimodal,westerhoff2025scam}. A note reading ``purple'' on a
red apple pits the image's own text against the object's pixels. Is this the same
arbitration problem, solved by the same heads? If the prior-override heads
implement a general ``answer from something other than the object's colour''
mechanism, they should carry the typographic override too. If they implement
routing from image to prompt text, as their attention shift suggests, they
should not, because here the competing signal is in the image.

We test this directly on Qwen2.5-VL-3B~\citep{bai2025qwen25vl} and
PaliGemma2-3B~\citep{steiner2024paligemma2}, holding the task, objects,
answer vocabulary, and
circuit-discovery pipeline fixed. We compare three answer-matched conflicts: P,
a prompted-prior conflict; T, a typographic conflict; and V, the image-varying
contrast from the appendix of \citet{lietzow2026}, which controls for T varying
the image. Our contributions are:
\begin{itemize}\itemsep0pt
\item An answer-matched P/T/V design and a pipeline that re-discovers
\citet{lietzow2026}'s circuits exactly in both models.
\item Circuit overlap read against a chance floor and a split-half ceiling fixed
before the T results existed: P and T are partly shared in both models, and the
image-varying circuit V lies inside T.
\item Cross-ablation against layer-matched random heads, showing that the shared
heads are necessary for P only in Qwen; in PaliGemma, the separation is partial.
\item A mechanism for Qwen's T-only heads: they read the note and write its
colour, and with no note do the same with the object's colour.
\end{itemize}
Every test, threshold and verdict phrase was pre-registered in a dated decision
log before the run it judges. The decision log will be released with the code.

\section{Related Work}
\paragraph{Perception--knowledge conflicts in VLMs.}
\citet{lietzow2026} localise prior override to sparse late heads in several VLMs,
and report that the heads they find under an image-varying contrast are largely
contained in the prompt-varying set. \citet{golovanevsky2025counterfact}
introduce the recoloured-object data we use. \citep{golovanevsky2024notice,hua2025,ortu2025,nooralahzadeh2026}
study related routing and arbitration mechanisms.\todo{read Hua, Ortu, Nooralahzadeh and
Golovanevsky 2024 in full; one accurate sentence each; check novelty against Hua
(image-vs-caption heads in Qwen2.5-VL)}

\paragraph{Typographic attacks.}
Text in an image can override object recognition in CLIP~\citep{goh2021multimodal}.
SCAM~\citep{westerhoff2025scam} evaluates real-world typographic attacks on
multimodal models, and \citet{dyslexify2025} locate typographic circuits in the
vision encoder.\todo{verify both} Our question concerns the language-model
side: which decoder heads carry the attack to the answer, and whether they are
the heads that carry prompted knowledge.

\paragraph{Circuit discovery and overlap.}
We use activation patching~\citep{meng2022rome,wang2023ioi,heimersheim2024patching}
at the level of attention-head outputs, following \citet{lietzow2026}.
\citet{nikankin2025} compare circuits across modalities by
intersection-over-union against random sets.\todo{verify} We add a split-half
ceiling, so a low overlap can be told apart from an unreliable circuit.

\section{Method}

\paragraph{Models and data.}
Qwen2.5-VL-3B-Instruct has \nLayers{} layers of \nHeadsPerLayer{} heads per
layer (\nHeads{} in total). PaliGemma2-3B uses checkpoint
\texttt{paligemma2-3b-mix-448}, with \pgLayers{} layers of \pgHeadsPerLayer{}
heads per layer (\pgHeads{} in total). Both models are loaded in 4-bit
NF4~\citep{dettmers2023qlora} with fp16 compute, as in \citet{lietzow2026}.
We use \nObjects{} objects from the colour split of
Visual-Counterfact~\citep{golovanevsky2025counterfact}, filtered as Lietzow et
al.'s loader filters it. \citet{lietzow2026} report 467 usable objects; we
could not reproduce their last two exclusions. Each object has an original
image showing its typical colour $\cco$ and a counterfactual image recoloured
to $\ccf$. The prompts are taken verbatim from their code: \emph{Visual}
``What color is this \{object\} here? Respond with the color name only.'' and
\emph{Prior} ``What color is \{a/an object\} usually? Respond with the color
name only.'' PaliGemma takes the prompt as plain text after the image, without
a chat template.

\paragraph{Three conflicts.}
Each conflict is a pair of runs that differ in one factor. In the \emph{source}
run the competing signal wins; in the \emph{target} run the image wins.
\begin{itemize}\itemsep0pt
\item \textbf{P (prompted prior)}~\citep{lietzow2026}: the counterfactual image
with the Prior prompt (source, answer $\cco$) versus the Visual prompt (target,
answer $\ccf$). We keep objects answered as requested under both prompts
(``correctly conflicting''): \nP.
\item \textbf{T (typographic)}: the original image with a post-it-style note
placed on the object (one fifth of the image width, centred on the object's
bounding box) under the Visual prompt. In the source run the note reads $\ccf$;
in the target run it reads a neutral, non-colour word of similar length. The
two runs differ only in the note's word. We keep objects where the neutral-note
run answers an accepted colour and the incongruent run answers the note's colour:
\nT.
\item \textbf{V (image-varying)}~\citep{lietzow2026}: the Visual prompt on the
original image (source, answer $\cco$) versus the counterfactual image (target,
answer $\ccf$), Lietzow et al.'s visual contrast. We keep objects whose original
image is answered with a real colour and whose counterfactual image is answered
with $\ccf$ (\nVselected{} on Qwen), and patch \nV{} of them, as many as T.
V's token roles match P's, not T's.
\end{itemize}
The answer tokens match between P and T but their roles swap: in P the source
run answers $\cco$ from memory, while in T the source run answers $\ccf$ from
the note. V is the control for T's image variation. If T shares heads with V
more than with P, the overlap follows where the signal enters, not the token
roles.

\paragraph{Circuit discovery.}
We port Lietzow et al.'s procedure. For each object, we cache every head's output
(the input slice of the attention output projection) at the last prompt token
in both runs. For every head, we patch the source run's value into the target
run and the target run's value into the source run, and record the logit
difference between the two runs' answer tokens. A head's restoration score in
each direction is the fraction of the source--target logit difference that the
patch moves. We average the directional scores over objects and take PC1 across
heads. Heads beyond $+2\sigma$ on PC1 are \emph{promoting}, and those below
$-2\sigma$ are \emph{suppressing}. With Qwen's causal attention, patching only
the last position cannot change earlier positions; we run the prompt prefix
once and recompute only the last token per patch, checking before each sweep
that this matches a full forward pass. PaliGemma's prefix attends in both
directions, so a last-token patch there reaches earlier positions in later
layers. We therefore use full forward passes for PaliGemma; its intervention is
not identical to Qwen's.

\paragraph{Overlap.}
Our primary statistic is the Jaccard index $J_{PT}$ of the promoting sets, with
a bootstrap confidence interval. Each bootstrap resample redraws objects
independently per condition and re-runs classification
(\nBoot{} resamples). We read $J_{PT}$ on a fixed scale.
The \emph{chance floor} is the overlap of random head sets of the observed sizes
(\nPerm{} draws). The \emph{ceiling} is the
overlap of the P circuit with itself when discovered on random halves of its
objects (\nHalfSplits{} splits). The verdict is one of four phrases, fixed in
advance: ``not distinguishable from chance at this $n$'', ``not distinguishable
from the P circuit's own self-agreement'', ``inconclusive at this $n$'', or
``above chance, below self-agreement''. Because P and T keep different objects,
we also compute $J_{PT}$ on the objects present in both. We read $J_{PV}$ and
$J_{VT}$ on the same scale. The V test was fixed in advance as
$\Delta J = J_{VT} - J_{PT}$, with a bootstrap confidence interval.

\paragraph{Cross-ablation.}
We zero-ablate a head set's outputs at the last token. Following Lietzow et al.,
an ablation \emph{defeats} the override when the source run's answer logit falls
below the target run's answer logit. Each test compares the observed rate with
\hTwoDraws{} random sets of the same size and layer distribution, drawn from
unclassified heads; the empirical p-value uses an add-one correction. The
preregistered primary tests ablate $\HP$ on T objects and $\HT$ on P
objects, at a pre-registered $\alpha=\hTwoAlpha$. With \hTwoDraws{} draws,
the smallest attainable p-value is \hTwoMinP. We also decompose
each set into shared, P-only and T-only heads. For PaliGemma,
we use an exact replay from the earliest ablated layer, checked against full
forward passes. Mean ablation with \pgMeanDraws{} random draws on PaliGemma is
the robustness check.

\paragraph{Mechanism.}
This analysis runs on Qwen only. For each head, at the last token, we measure
(i) the attention mass of the last-token query on \emph{note} tokens, (ii) on
\emph{object} tokens, and (iii) its \emph{direct effect} on the answer. Note
tokens are image tokens at least half covered by the note on the model's own
token grid. Object tokens are image tokens at least half covered by recoloured
pixels and untouched by the note. The direct effect is a logit lens
~\citep{nostalgebraist2020logitlens}: the head's contribution to the residual
stream, passed through the final norm and the unembedding, read as the logit of
the source answer minus the target answer. Attention weights require eager
attention. Before measuring, we check that the weights rebuild every head's
output and that the logit lens reproduces the model's own logits. Statistics
and random sets are those of the cross-ablation.
\section{Results}
\subsection{Replication of the prior-override circuit}\label{sec:replication}
On the \nP{} correctly conflicting objects (which \citet{lietzow2026} report
as 73 for this model), patching over all \nHeads{} heads finds \nPromP{} promoting and \nSupP{}
suppressing heads, \promMatchesLietzow{} their published lists. Ablating the
promoting group defeats the prior override on \repPriorFlip\% of objects and
the visual answer on \repVisualFlip\% (theirs: \lietzowPriorFlip\% and
\lietzowVisualFlip\%). The strongest single heads are \topSingleOne{}
(\topSingleOneFlip\%; theirs \topSingleOneLietzow\%) and \topSingleTwo{}
(\topSingleTwoFlip\%; theirs \topSingleTwoLietzow\%). We use our re-discovered
set as $\HP$, so that P and T come from an identical pipeline.

On the \pgNP{} correctly conflicting objects (which \citet{lietzow2026} also
report as 121),
patching over all \pgHeads{} heads finds \pgNPromP{} promoting and
\pgNSupP{} suppressing heads, \pgMatchesLietzow{} their published lists. Group ablation defeats the prior override on \pgPriorFlip\%
and the visual answer on \pgVisualFlip\% (theirs: \pgLietzowPriorFlip\% and
\pgLietzowVisualFlip\%). The strongest single heads are \pgTopOne{}
(\pgTopOneFlip\%; theirs \pgTopOneLietzow\%) and \pgTopTwo{}
(\pgTopTwoFlip\%; theirs \pgTopTwoLietzow\%).

\begin{table*}[t]
\centering
\small
\begin{tabular}{lcc}
\toprule
 & Qwen2.5-VL-3B & PaliGemma2-3B \\
\midrule
P objects & \nP & \pgNP \\
P promoting heads (lists identical to Lietzow) & \nPromP & \pgNPromP \\
Prior / visual flip, \% & \repPriorFlip{} / \repVisualFlip & \pgPriorFlip{} / \pgVisualFlip \\
P split-half ceiling, 2.5th pct & \ceilingLo & \pgCeilingLo \\
T objects & \nT & \pgNT \\
T promoting heads & \nPromT & \pgNPromT \\
Shared heads & \nShared & \pgNShared \\
$J_{PT}$ [95\% CI] & \JPT{} [\JPTLo, \JPTHi] & \pgJPT{} [\pgJPTLo, \pgJPTHi] \\
Chance, 95th pct & \chanceEdge & \pgChanceEdge \\
$\HP$ on T vs random, \% & \hTwoPtoTObs{} vs \hTwoPtoTNull & \pgPtoTObs{} vs \pgPtoTNull \\
$\HT$ on P vs random, \% & \hTwoTtoPObs{} vs \hTwoTtoPNull & \pgTtoPObs{} vs \pgTtoPNull \\
\bottomrule
\end{tabular}
\caption{Replication, circuit overlap and cross-ablation summary for both models.}
\label{tab:replication}
\end{table*}

\subsection{Notes fool the model}
A behavioural screen of all \nObjects{} objects shows that the note works by its
content, not its presence. An incongruent note moves the logit difference toward
the note's colour by \attackShift{} logits on average (minimum \attackShiftMin),
while a note naming the object's own colour moves it \congruentShift{} the other
way (maximum \congruentShiftMax). Of the \nCorrectNeutral{} objects answered
correctly with a neutral note in at least one of three placements,
\screenFoolRate\% are fooled by the incongruent note in at least one such
placement. Placement matters more than size: small and medium notes on the
object fool \screenFoolSmall\% and \screenFoolMedium\% of eligible objects, a
note beside it \screenFoolBeside\%. Paired within objects, \pairedOnOnly{} are
fooled only on the object and \pairedBesideOnly{} only beside it. The behavioural
screen ran with bf16 compute; every circuit experiment used fp16. \nT{} objects
meet the T selection rule for discovery. On PaliGemma, \pgNT{} objects meet the
same T selection rule with a small note on the object.

\subsection{The circuits overlap partly}\label{sec:overlap}
Discovery on T finds \nPromT{} promoting heads. \nShared{} of them are also P
heads, including the two strongest single P heads (Fig.~\ref{fig:heads}). The
P-only heads lie in layers \layersPonly{} and the T-only heads in layers
\layersTonly{}.
$J_{PT}=\JPT$ [\JPTLo, \JPTHi]. This lies far from both ends of the scale
(Fig.~\ref{fig:overlap}): the chance floor's 95th percentile is \chanceEdge{}
and the P ceiling's 2.5th percentile is \ceilingLo. The pre-registered verdict
is \emph{``\JPTVerdict''}. Neither pure hypothesis fits: routing-specific heads
predict overlap near chance, and a general override mechanism predicts overlap
near the ceiling. Object selection does not drive the result: on the
\nSharedObjects{} objects in both sets, $J_{PT}=\JPTshared$ [\JPTsharedLo,
\JPTsharedHi]. Nor is T the noisy side, since T's circuit agrees with itself
across halves at median $J=\TselfJ$. Across all heads, the rank correlation of
PC1 scores is \rankRho{} (P's own, full length: \rankCeiling).

\begin{figure}[t]
\centering
\figplaceholder{0.85\columnwidth}{6cm}{fig\_heads: PC1 score of every head, P vs T}
\caption{Every head's PC1 score under P and under T, in units of each
condition's standard deviation; dashed lines are the $\pm2\sigma$ cuts. The
circuits share \nShared{} heads (labelled). The P-only and T-only heads sit just
above the cut in one condition and near zero in the other.}
\label{fig:heads}
\end{figure}

\begin{figure}[t]
\centering
\figplaceholder{0.85\columnwidth}{4.5cm}{fig\_overlap: J with chance floor and split-half ceiling}
\caption{$J_{PT}$ with bootstrap confidence intervals, against the chance floor
(random sets of the observed sizes, with a bar to the upper percentile and a
tick at the mean) and the P circuit's split-half self-agreement (percentile
interval with a tick at the median). Hollow markers show objects kept by both
conditions. The V rows will also be included.}
\label{fig:overlap}
\end{figure}

On PaliGemma, $J_{PT}=\pgJPT$ [\pgJPTLo, \pgJPTHi], the chance 95th percentile
is \pgChanceEdge{} and the ceiling 2.5th percentile is \pgCeilingLo. The
pre-registered verdict is \emph{``\pgJPTVerdict''}. On the \pgNSharedObjects{}
objects in both sets, $J_{PT}=\pgJPTshared$ [\pgJPTsharedLo, \pgJPTsharedHi].
Unlike on Qwen, P's strongest single head, \pgTopOne, is not shared. Head
indices do not correspond across models, and the two $J_{PT}$ values were not
tested against each other.

On Qwen, V finds \nPromV{} promoting heads, all in layers \layersV, and every
one is also a T head. $J_{VT}=\JVT$ [\JVTLo, \JVTHi] and
$J_{PV}=\JPV$ [\JPVLo, \JPVHi], with pre-registered verdicts
\emph{``\JVTVerdict''} and \emph{``\JPVVerdict''}. The pre-registered test
gives $\Delta J=\dJ$ [\dJLo, \dJHi], verdict \emph{``\dJVerdict''}.
\nTonlyInV{} of the \nTonly{} T-only heads and \nSharedInV{} of the \nShared{}
shared heads are V heads. Because V's token roles match P's, a token-role effect
would have pushed the other way. $J_{VT}$ sits at the top of its interval because
it is the largest $J$ those set sizes allow. No P--V versus P--T comparison was
pre-registered, so we claim nothing about $J_{PV}$ versus $J_{PT}$. As a
reference only, V's containment in P is \containVinP{}, while \citet{lietzow2026}
report 61\% (whether they pool head classes is unclear).
\subsection{Shared selection, separate causes}\label{sec:h2}
On Qwen, both primary cross-ablation tests are null (Fig.~\ref{fig:h2}).
Ablating $\HP$ defeats the typographic override on \hTwoPtoTObs\% of T
objects [\hTwoPtoTLo, \hTwoPtoTHi] (random: \hTwoPtoTNull\%,
$p=\hTwoPtoTP$). Ablating $\HT$ defeats the prior override on
\hTwoTtoPObs\% of P objects [\hTwoTtoPLo, \hTwoTtoPHi] (random:
\hTwoTtoPNull\%, $p=\hTwoTtoPP$). Both are \emph{``not distinguishable
from layer-matched random heads at this $n$''}. Each set does defeat its own
override: \hTwoPtoPObs\% and \hTwoTtoTObs\%. \hTwoMeanSentence

The decomposition explains the nulls. The \nShared{} shared heads alone
defeat the prior override on \hTwoStoPObs\% of P objects, but the typographic
override on \hTwoStoTObs\% of T objects (\hTwoStoTCount{} of \hTwoStoTN).
The T override runs through the \nTonly{} T-only heads, which alone defeat
it on \hTwoTOtoTObs\%. The P-only heads defeat it on \hTwoPOtoTObs\%
($p=\hTwoPOtoTP$), which does not reach $\alpha$. $\HT$ on P reads null for
a different reason than $\HP$ on T. The T-only heads push P's source runs
\emph{toward} the image's colour (mean logit shift \hTwoTOonPShift, against
\hTwoStoPShift{} for the shared heads), so ablating all of $\HT$ cancels the
shared heads' effect.

On PaliGemma, ablating $\HP$ on T defeats the override on \pgPtoTObs\%
[\pgPtoTLo, \pgPtoTHi] versus random \pgPtoTNull\%, $p=\pgPtoTP$:
``\pgPtoTVerdict''. Ablating $\HT$ on P defeats the override on \pgTtoPObs\%
[\pgTtoPLo, \pgTtoPHi] versus random \pgTtoPNull\%, $p=\pgTtoPP$:
``\pgTtoPVerdict'', unlike Qwen. The T-only heads defeat the P override on
\pgTOtoPCount{} of \pgTOtoPN{} objects, so $\HT$ has no separate route to P.
As on Qwen, they push P's answer toward the image's colour (mean logit shift
\pgTOonPShift{} against \pgStoPShift{} for the shared heads), but here this
cancels the shared heads' effect only in part (shared heads alone:
\pgStoPObs\%). The shared heads carry a little of T on PaliGemma
(\pgStoTObs\% vs random \pgStoTNull\%); on Qwen they carried none. Under
mean ablation, $\HP$ on T gives \pgMeanPtoTObs\% vs \pgMeanPtoTNull\% and
$\HT$ on P gives \pgMeanTtoPObs\% vs \pgMeanTtoPNull\%
($p=\pgMeanPtoTP$, the smallest attainable with \pgMeanDraws{} draws).
\pgMeanSentence{}

\begin{figure*}[t]
\centering
\figplaceholder{0.85\textwidth}{5.5cm}{fig\_cross\_ablation: share of overrides defeated per head set vs random sets}
\caption{Cross-ablation. Share of objects whose override is defeated by
zero-ablating each head set, with 95\% bootstrap CIs. Grey: \hTwoDraws{}
layer-matched random sets (bar to the 95th percentile, tick at the mean). Hollow
markers have no random control (descriptive).}
\label{fig:h2}
\end{figure*}

\subsection{The T-only heads move colour from the image}\label{sec:h3}
This mechanism analysis ran on Qwen only. On \hThreeNT{} T objects
(\hThreeDropT{} dropped because their answers changed under eager attention or
became non-finite), both primary mechanism tests pass (Fig.~\ref{fig:h3}a,b).
The T-only heads place \hThreeAObs{} of their last-token attention on note
tokens [\hThreeALo, \hThreeAHi], against \hThreeANull{} for random heads
($p=\hThreeAP$). Their direct effect toward the note's colour is
\hThreeWObs{} logits [\hThreeWLo, \hThreeWHi], against \hThreeWNull{}.
The pre-registered mechanism sentence is \emph{``\hThreeMechanism''}. The
note makes up \noteShare\% of the image tokens, but receives \enrichInc{}
times its share of the T-only heads' image attention with a colour word and
\enrichNeu{} times with a neutral word. The heads read text in the image
regardless of its content, and read it more when it names a colour (rise
\hThreeRiseObs, random \hThreeRiseNull).

The shared heads neither attend to the note (\hThreeASObs, random
\hThreeASNull, $p=\hThreeASP$) nor write its colour (\hThreeWSObs{} logits),
consistent with their null ablation on T. The P-only heads attend to the note
(\hThreeAPOObs, random \hThreeAPONull) but do not write its colour
(\hThreeWPOObs).

With no note, the T-only heads do the same thing with the object's own colour
(Fig.~\ref{fig:h3}c). On the clean original image they attend to the recoloured
object region (\hThreeObjCleanObs, random \hThreeObjCleanNull) and write the
image's colour (\hThreeColCleanObs{} logits, random \hThreeColCleanNull). On
P's Visual-prompt runs they do likewise (\hThreeObjPVObs; \hThreeColPVObs{}
logits). As a check on the attention measurement, $\HP$ attends to the image
more under the Visual than the Prior prompt, the direction \citet{lietzow2026}
report (difference \hThreeRoutingObs, random \hThreeRoutingNull,
$p=\hThreeRoutingP$). V gives independent support: \nTonlyInV{} of these
\nTonly{} heads are selected by a contrast with no note at all.

\begin{figure*}[t]
\centering
\figplaceholder{0.85\textwidth}{5cm}{fig\_h3: attention on note tokens and direct effect, per head set}
\caption{Mechanism. (a) Last-token attention mass on note tokens, mean over the
set's heads; (b) the set's summed direct effect toward the note's colour, both on
T objects with the incongruent note. (c) The T-only heads' direct effect with the
note, and with no note on T's clean images and P's Visual-prompt runs (negative =
toward the image's colour). Grey as in Fig.~\ref{fig:h2}; whiskers are 95\%
bootstrap CIs.}
\label{fig:h3}
\end{figure*}

\section{Discussion}

\paragraph{Selection is not cause.}
On Qwen, patching selected the shared heads in both P and T, yet ablating them
affects only P, and under T they neither attend to the note nor write its
colour. Patching asks whether a head's output \emph{carries} the difference
between two runs; ablation asks whether the answer \emph{needs} it. For the
shared heads under T, the two disagree. V sharpens this point: the plain image
contrast does not select two of the shared heads, L26H0 and L26H5 (P's
strongest single head), although T does. Why patching selects them under T
remains open. Path patching from these heads to the T-only heads would separate
redundancy from an indirect route. Overlap statistics computed from patching
alone would have overstated how much the two overrides share.

\paragraph{Two routes to an override.}
The prior-override heads let prompted knowledge beat the image. The typographic
attack does not use them. It uses heads that, by our measurements, move colour
from wherever it appears in the image to the answer. A note is effective
because it hands those heads a colour, not because it triggers the
prior-override machinery. This fits the attack's dependence on placement: a
note on the object fools more often than one beside it. The V circuit supports
this reading: heads that carry an object's real colour from the pixels are a
subset of those that carry a note's colour. Yet ablating the T-only heads barely
affects ordinary answers (clean accuracy kept at \hTwoCleanTO\%, random
\hTwoCleanRandom\%), so visual colour must also reach the answer by other
routes. This is consistent with visual grounding being the default.

\paragraph{A second model.}
On PaliGemma, the overlap result repeats the same verdict, ``\pgJPTVerdict'',
but cross-ablation separation is only partial and not robust to ablation type,
so the causal account rests on Qwen. \citet{lietzow2026} find that PaliGemma
routes differently; our PaliGemma intervention also differs because its prefix
is bidirectional.

\paragraph{Limitations.}
Both models are quantised. Interventions and attention are measured at the last
token only. P and T share the colour pair but swap its roles, so heads that
\emph{write} an answer may differ partly for that reason. The image-varying
contrast V, whose roles match P, addresses this, but the mechanism and V
analyses ran on Qwen only. Head indices do not correspond across models, and
the PaliGemma cross-ablation verdict depends on the ablation type. Head selection
makes a positive direct effect for the T-only heads partly expected; the
attention result does not share this problem. In the mechanism analysis,
``object'' tokens cover only the recoloured region and can omit parts of an
object left unchanged. Many tests sit at the smallest p-value \hTwoDraws{}
random sets allow, so the excess over random sets and its interval are the
effect sizes.

\section{Conclusion}
Prompted priors and in-image text both override what a VLM sees, and patching
selects partly the same heads for both. In Qwen, the prior-override heads are
not needed for the typographic attack, which instead runs through heads that
carry colour from the image, whether painted on an object or written on a note.
In PaliGemma, the overlap repeats the verdict ``\pgJPTVerdict'', but causal
separation is only partial.

\bibliography{refs}
\end{document}